% ===================================================================
% 统一排版样式指南 - Unified Style Guide
% ===================================================================
% 本文件定义论文中所有图表的统一样式规范
% 用于确保排版的一致性和美观性

% ===================================================================
% 1. 表格统一样式 (Table Unified Style)
% ===================================================================

% 表格标题字体：9pt，粗体标签
\DeclareCaptionFont{unifiedtablefont}{\fontsize{9pt}{11.5pt}\selectfont}
\captionsetup[table]{%
  font=unifiedtablefont,
  labelfont=bf,
  textfont=normalfont,
  justification=raggedright,
  singlelinecheck=false,
  width=\linewidth,
  skip=5pt,
  labelformat=tablelabel,
  labelsep=none,
  format=tabletwoline
}

% 表格主体样式命令
\renewcommand{\TableStyle}{%
  \centering
  \setstretch{1}%
  \fontsize{9pt}{11.5pt}\selectfont
  \setlength{\parindent}{0pt}%
  \setlength{\parskip}{0pt}%
  \setlength{\tabcolsep}{4pt}%
  \renewcommand{\arraystretch}{1.12}%
  \setlength{\heavyrulewidth}{0.75pt}%
  \setlength{\lightrulewidth}{0.45pt}%
  \setlength{\cmidrulewidth}{0.45pt}%
  \setlength{\aboverulesep}{0.6ex}%
  \setlength{\belowrulesep}{0.6ex}%
}

% ===================================================================
% 2. 图片统一样式 (Figure Unified Style)
% ===================================================================

% 图片标题字体：9pt
\DeclareCaptionFont{unifiedfigurefont}{\fontsize{9pt}{11.5pt}\selectfont}
\captionsetup[figure]{%
  skip=4.5pt,
  justification=centering,
  singlelinecheck=false,
  font=unifiedfigurefont,
  labelfont=normalfont,
  textfont=normalfont,
  labelformat=figlabel,
  labelsep=period
}

% 单张图片样式命令
\newcommand{\SingleFigure}[3]{%
  % #1: 图片文件路径
  % #2: 图片宽度（默认0.85\textwidth）
  % #3: caption文本
  \begin{figure}[htbp]
  \centering
  \includegraphics[width=#2]{#1}
  \caption{#3}
  \end{figure}
}

% 并排图片样式命令（2列）
\newcommand{\DualFigure}[6]{%
  % #1: 左图路径, #2: 左图标题
  % #3: 右图路径, #4: 右图标题
  % #5: 整体caption, #6: label
  \begin{figure}[htbp]
  \centering
  \begin{minipage}[t]{0.48\textwidth}
    \centering
    {\fontsize{9.5}{11.5}\selectfont\bfseries #2}\\[0.2em]
    \includegraphics[width=\linewidth]{#1}
  \end{minipage}%
  \hfill
  \begin{minipage}[t]{0.48\textwidth}
    \centering
    {\fontsize{9.5}{11.5}\selectfont\bfseries #4}\\[0.2em]
    \includegraphics[width=\linewidth]{#3}
  \end{minipage}
  \caption{#5}
  \label{#6}
  \end{figure}
}

% 对比实验图专用样式（3列×2行，紧凑无留白）
\newcommand{\ComparisonFigureGrid}[2]{%
  % #1: minipage内容（已定义的子图结构）
  % #2: 总caption
  \begin{figure}[htbp]
  \centering
  \makebox[\textwidth][c]{%
    \begin{minipage}{18.5cm}
    \centering
    #1
    \end{minipage}%
  }
  \caption{#2}
  \end{figure}
}

% 子图标题统一格式
\newcommand{\SubFigTitle}[1]{%
  {\fontsize{10}{12}\selectfont\bfseries #1}\\[-0.1em]%
}

% ===================================================================
% 3. 浮动体间距控制 (Float Spacing Control)
% ===================================================================

% 文本与浮动体之间的垂直间距
\setlength{\textfloatsep}{9pt plus 2pt minus 1pt}

% 连续浮动体之间的间距
\setlength{\floatsep}{8pt plus 2pt minus 1pt}

% 顶部/底部浮动体与文本的间距
\setlength{\intextsep}{8pt plus 2pt minus 1pt}

% ===================================================================
% 4. 图表编号格式 (Numbering Format)
% ===================================================================

% 图片编号格式：Fig. X
\DeclareCaptionLabelFormat{figlabel}{Fig.\ #2}
\crefname{figure}{Fig.}{Figs.}
\Crefname{figure}{Fig.}{Figs.}

% 表格编号格式：Table X
\DeclareCaptionLabelFormat{tablelabel}{Table\ #2}
\crefname{table}{Table}{Tables}
\Crefname{table}{Table}{Tables}

% ===================================================================
% 5. 特殊说明
% ===================================================================
% - 所有表格必须使用\TableStyle命令
% - 所有对比实验图（多子图网格）使用\ComparisonFigureGrid
% - 子图标题使用\SubFigTitle命令确保字号一致
% - 禁止手动调整个别图表的字号，必须使用统一命令
% ===================================================================
